% Options for packages loaded elsewhere
% Options for packages loaded elsewhere
\PassOptionsToPackage{unicode}{hyperref}
\PassOptionsToPackage{hyphens}{url}
%
\documentclass[
  letterpaper,
  DIV=11,
  numbers=noendperiod]{scrartcl}
\usepackage{xcolor}
\usepackage[top=1in,bottom=1in,left=1in,right=1in]{geometry}
\usepackage{amsmath,amssymb}
\setcounter{secnumdepth}{5}
\usepackage{iftex}
\ifPDFTeX
  \usepackage[T1]{fontenc}
  \usepackage[utf8]{inputenc}
  \usepackage{textcomp} % provide euro and other symbols
\else % if luatex or xetex
  \usepackage{unicode-math} % this also loads fontspec
  \defaultfontfeatures{Scale=MatchLowercase}
  \defaultfontfeatures[\rmfamily]{Ligatures=TeX,Scale=1}
\fi
\usepackage{lmodern}
\ifPDFTeX\else
  % xetex/luatex font selection
  \setmainfont[]{TeX Gyre Pagella}
  \setmonofont[]{Courier New}
\fi
% Use upquote if available, for straight quotes in verbatim environments
\IfFileExists{upquote.sty}{\usepackage{upquote}}{}
\IfFileExists{microtype.sty}{% use microtype if available
  \usepackage[]{microtype}
  \UseMicrotypeSet[protrusion]{basicmath} % disable protrusion for tt fonts
}{}
\makeatletter
\@ifundefined{KOMAClassName}{% if non-KOMA class
  \IfFileExists{parskip.sty}{%
    \usepackage{parskip}
  }{% else
    \setlength{\parindent}{0pt}
    \setlength{\parskip}{6pt plus 2pt minus 1pt}}
}{% if KOMA class
  \KOMAoptions{parskip=half}}
\makeatother
% Make \paragraph and \subparagraph free-standing
\makeatletter
\ifx\paragraph\undefined\else
  \let\oldparagraph\paragraph
  \renewcommand{\paragraph}{
    \@ifstar
      \xxxParagraphStar
      \xxxParagraphNoStar
  }
  \newcommand{\xxxParagraphStar}[1]{\oldparagraph*{#1}\mbox{}}
  \newcommand{\xxxParagraphNoStar}[1]{\oldparagraph{#1}\mbox{}}
\fi
\ifx\subparagraph\undefined\else
  \let\oldsubparagraph\subparagraph
  \renewcommand{\subparagraph}{
    \@ifstar
      \xxxSubParagraphStar
      \xxxSubParagraphNoStar
  }
  \newcommand{\xxxSubParagraphStar}[1]{\oldsubparagraph*{#1}\mbox{}}
  \newcommand{\xxxSubParagraphNoStar}[1]{\oldsubparagraph{#1}\mbox{}}
\fi
\makeatother


\usepackage{longtable,booktabs,array}
\usepackage{calc} % for calculating minipage widths
% Correct order of tables after \paragraph or \subparagraph
\usepackage{etoolbox}
\makeatletter
\patchcmd\longtable{\par}{\if@noskipsec\mbox{}\fi\par}{}{}
\makeatother
% Allow footnotes in longtable head/foot
\IfFileExists{footnotehyper.sty}{\usepackage{footnotehyper}}{\usepackage{footnote}}
\makesavenoteenv{longtable}
\usepackage{graphicx}
\makeatletter
\newsavebox\pandoc@box
\newcommand*\pandocbounded[1]{% scales image to fit in text height/width
  \sbox\pandoc@box{#1}%
  \Gscale@div\@tempa{\textheight}{\dimexpr\ht\pandoc@box+\dp\pandoc@box\relax}%
  \Gscale@div\@tempb{\linewidth}{\wd\pandoc@box}%
  \ifdim\@tempb\p@<\@tempa\p@\let\@tempa\@tempb\fi% select the smaller of both
  \ifdim\@tempa\p@<\p@\scalebox{\@tempa}{\usebox\pandoc@box}%
  \else\usebox{\pandoc@box}%
  \fi%
}
% Set default figure placement to htbp
\def\fps@figure{htbp}
\makeatother





\setlength{\emergencystretch}{3em} % prevent overfull lines

\providecommand{\tightlist}{%
  \setlength{\itemsep}{0pt}\setlength{\parskip}{0pt}}



 


% --- Figure Spacing (Your existing settings) ---
\setlength{\intextsep}{12pt} 
\setlength{\textfloatsep}{14pt}
\setlength{\floatsep}{12pt}

% --- Headers and Footers ---
\usepackage[headsepline, footsepline]{scrlayer-scrpage}
\usepackage{cancel}
\clearpairofpagestyles
\ihead{CBE 20260}  % Inner Head (Left)
\ohead{Spring 2026}   % Outer Head (Right)
\cfoot{\pagemark}     % Center Foot (Page Number)
\pagestyle{scrheadings}

% --- Optional: Make links look nicer ---
\usepackage{hyperref}
\hypersetup{
    colorlinks=true,
    linkcolor=blue,
    filecolor=magenta,      
    urlcolor=blue,
}
\usepackage{fvextra}
\DefineVerbatimEnvironment{Highlighting}{Verbatim}{breaklines,commandchars=\\\{\}}
\KOMAoption{captions}{tableheading}
\makeatletter
\@ifpackageloaded{tcolorbox}{}{\usepackage[skins,breakable]{tcolorbox}}
\@ifpackageloaded{fontawesome5}{}{\usepackage{fontawesome5}}
\definecolor{quarto-callout-color}{HTML}{909090}
\definecolor{quarto-callout-note-color}{HTML}{0758E5}
\definecolor{quarto-callout-important-color}{HTML}{CC1914}
\definecolor{quarto-callout-warning-color}{HTML}{EB9113}
\definecolor{quarto-callout-tip-color}{HTML}{00A047}
\definecolor{quarto-callout-caution-color}{HTML}{FC5300}
\definecolor{quarto-callout-color-frame}{HTML}{acacac}
\definecolor{quarto-callout-note-color-frame}{HTML}{4582ec}
\definecolor{quarto-callout-important-color-frame}{HTML}{d9534f}
\definecolor{quarto-callout-warning-color-frame}{HTML}{f0ad4e}
\definecolor{quarto-callout-tip-color-frame}{HTML}{02b875}
\definecolor{quarto-callout-caution-color-frame}{HTML}{fd7e14}
\makeatother
\makeatletter
\@ifpackageloaded{caption}{}{\usepackage{caption}}
\AtBeginDocument{%
\ifdefined\contentsname
  \renewcommand*\contentsname{Table of contents}
\else
  \newcommand\contentsname{Table of contents}
\fi
\ifdefined\listfigurename
  \renewcommand*\listfigurename{List of Figures}
\else
  \newcommand\listfigurename{List of Figures}
\fi
\ifdefined\listtablename
  \renewcommand*\listtablename{List of Tables}
\else
  \newcommand\listtablename{List of Tables}
\fi
\ifdefined\figurename
  \renewcommand*\figurename{Fig.}
\else
  \newcommand\figurename{Fig.}
\fi
\ifdefined\tablename
  \renewcommand*\tablename{Table}
\else
  \newcommand\tablename{Table}
\fi
}
\@ifpackageloaded{float}{}{\usepackage{float}}
\floatstyle{ruled}
\@ifundefined{c@chapter}{\newfloat{codelisting}{h}{lop}}{\newfloat{codelisting}{h}{lop}[chapter]}
\floatname{codelisting}{Listing}
\newcommand*\listoflistings{\listof{codelisting}{List of Listings}}
\makeatother
\makeatletter
\makeatother
\makeatletter
\@ifpackageloaded{caption}{}{\usepackage{caption}}
\@ifpackageloaded{subcaption}{}{\usepackage{subcaption}}
\makeatother
\usepackage{bookmark}
\IfFileExists{xurl.sty}{\usepackage{xurl}}{} % add URL line breaks if available
\urlstyle{same}
\hypersetup{
  pdftitle={Lecture 33: Chemical Potential, Gibbs Energy, and the Clapeyron Equation},
  pdfauthor={Edward Maginn},
  hidelinks,
  pdfcreator={LaTeX via pandoc}}


\title{Lecture 33: Chemical Potential, Gibbs Energy, and the Clapeyron
Equation}
\usepackage{etoolbox}
\makeatletter
\providecommand{\subtitle}[1]{% add subtitle to \maketitle
  \apptocmd{\@title}{\par {\large #1 \par}}{}{}
}
\makeatother
\subtitle{CBE 20260 / Thermodynamics I}
\author{Edward Maginn}
\date{}
\begin{document}
\maketitle

\renewcommand*\contentsname{Table of contents}
{
\setcounter{tocdepth}{3}
\tableofcontents
}

\section{Review and Continuation}\label{review-and-continuation}

Last lecture, we established two key ideas about phase equilibrium:

\begin{enumerate}
\def\labelenumi{\arabic{enumi}.}
\item
  For two phases in equilibrium, temperature, pressure, and chemical
  potential are equal: \[
  T^\alpha = T^\beta,\qquad
  P^\alpha = P^\beta,\qquad
  \mu^\alpha = \mu^\beta
  \]
\item
  For a pure fluid, the chemical potential is just the molar Gibbs
  energy: \[
  \mu = \underline{G}
  \]
\end{enumerate}

This lecture picks up from there. We will:

\begin{itemize}
\tightlist
\item
  connect chemical potential to all of the major thermodynamic
  potentials,
\item
  write the general fundamental relations when (N) is allowed to vary,
\item
  show that maximizing entropy for a closed system at constant (T) and
  (P) implies minimization of (G),
\item
  derive a very important result relating the \textbf{rate of change of
  Gibbs energy} to the \textbf{rate of entropy generation},
\item
  derive the \textbf{Clapeyron equation}, and
\item
  simplify it to the \textbf{Clausius--Clapeyron equation} for
  vapor-liquid equilibrium.
\end{itemize}

\begin{center}\rule{0.5\linewidth}{0.5pt}\end{center}

\section{Chemical Potential and Thermodynamic
Potentials}\label{chemical-potential-and-thermodynamic-potentials}

For a pure fluid,

\[
\mu = \underline{G}
\]

More generally, the chemical potential is the \textbf{partial molar
Gibbs energy}:

\begin{equation}\phantomsection\label{eq-chemical-potential}{
\mu = \left( \frac{\partial G}{\partial N} \right)_{T,P} 
}\end{equation}

But \(\mu\) also appears naturally when we differentiate the other
thermodynamic potentials with respect to the number of moles, holding
their natural variables constant:

\subsection{\texorpdfstring{Chemical potential from \(U\), \(H\), \(A\),
and
\(G\)}{Chemical potential from U, H, A, and G}}\label{chemical-potential-from-u-h-a-and-g}

For a one-component system,

\begin{equation}\phantomsection\label{eq-mu-from-U}{
\mu = \left( \frac{\partial U}{\partial N} \right)_{S,V} 
}\end{equation}

\begin{equation}\phantomsection\label{eq-mu-from-H}{
\mu = \left( \frac{\partial H}{\partial N} \right)_{S,P}   
}\end{equation}

\begin{equation}\phantomsection\label{eq-mu-from-A}{
\mu = \left( \frac{\partial A}{\partial N} \right)_{T,V} 
}\end{equation}

\begin{equation}\phantomsection\label{eq-mu-from-G}{
\mu = \left( \frac{\partial G}{\partial N} \right)_{T,P} 
}\end{equation}

So \(\mu\) is the conjugate variable to \(N\), just as:

\begin{itemize}
\tightlist
\item
  \(T\) is conjugate to \(S\),
\item
  \(-P\) is conjugate to \(V\).
\end{itemize}

This is why chemical potential is the natural thermodynamic quantity
governing transfer of matter. When there is no net matter transferred
between two phases, the chemical potentials must be equal.

\begin{center}\rule{0.5\linewidth}{0.5pt}\end{center}

\section{General Fundamental Relations with Chemical
Potential}\label{general-fundamental-relations-with-chemical-potential}

We first learned the fundamental relation for a closed, one-component
system with fixed composition:

\[
d\underline{U} = T\,d\underline{S} - P\,d\underline{V}
\]

If the number of moles is allowed to vary, then we must add a third term
involving (dN). The full differential relations for the extensive
potentials are:

\begin{equation}\phantomsection\label{eq-fundamental-U}{
dU = T\,dS - P\,dV + \mu\,dN
}\end{equation}

\begin{equation}\phantomsection\label{eq-fundamental-H}{
dH = T\,dS + V\,dP + \mu\,dN
}\end{equation}

\begin{equation}\phantomsection\label{eq-fundamental-A}{
dA = -S\,dT - P\,dV + \mu\,dN
}\end{equation}

\begin{equation}\phantomsection\label{eq-fundamental-G}{
dG = -S\,dT + V\,dP + \mu\,dN
}\end{equation}

These are the general fundamental relations for a one-component system.

\subsection{Important special case}\label{important-special-case}

For a \textbf{closed system}, (\(dN=0\)), these reduce to the familiar
forms:

\[
dU = T\,dS - P\,dV
\]

\[
dH = T\,dS + V\,dP
\]

\[
dA = -S\,dT - P\,dV
\]

\[
dG = -S\,dT + V\,dP
\]

\begin{center}\rule{0.5\linewidth}{0.5pt}\end{center}

\section{Phase Equilibrium in Terms of Chemical
Potential}\label{phase-equilibrium-in-terms-of-chemical-potential}

The general condition for phase equilibrium between any two phases
\(\alpha\) and \(\beta\) is:

\[
\mu^\alpha = \mu^\beta
\]

For a pure fluid, since \$\mu = \underline{G}, this can also be written
as

\[
\underline{G}^\alpha = \underline{G}^\beta
\]

Thus:

\begin{itemize}
\item
  vapor-liquid equilibrium: \[
  \underline{G}^v = \underline{G}^l
  \]
\item
  solid-liquid equilibrium: \[
  \underline{G}^s = \underline{G}^l
  \]
\item
  solid-vapor equilibrium: \[
  \underline{G}^s = \underline{G}^v
  \]
\item
  triple point: \[
  \underline{G}^s = \underline{G}^l = \underline{G}^v
  \]
\end{itemize}

If the molar Gibbs energies of the phases are not equal at a given \(T\)
and \(P\), then the system is not at phase equilibrium.

\begin{center}\rule{0.5\linewidth}{0.5pt}\end{center}

\section{Gibbs Energy and Spontaneous
Processes}\label{gibbs-energy-and-spontaneous-processes}

Now we want to connect phase equilibrium to the \textbf{direction of
spontaneous change}.

The key question is:

\begin{quote}
At constant \(T\) and \(P\), what thermodynamic potential determines
whether a process will proceed spontaneously?
\end{quote}

We will show that the answer is the Gibbs energy.

\begin{center}\rule{0.5\linewidth}{0.5pt}\end{center}

\section{Closed System Energy and Entropy
Balances}\label{closed-system-energy-and-entropy-balances}

Consider a \textbf{closed system} at constant temperature and pressure.
Ignore kinetic and potential energy changes.

\subsection{Energy balance}\label{energy-balance}

For a closed system,

\[
\frac{dU}{dt} = \dot{Q} + \dot{W}
\]

For a simple compressible system, if the only work is boundary (``PV'')
work,

\[
\dot{W} = -P\frac{dV}{dt}
\]

so

\begin{equation}\phantomsection\label{eq-energy-balance}{
\frac{dU}{dt} = \dot{Q} - P\frac{dV}{dt}
}\end{equation}

\subsection{Entropy balance}\label{entropy-balance}

The entropy balance for a closed system is

\begin{equation}\phantomsection\label{eq-entropy-balance}{
\frac{dS}{dt} = \frac{\dot{Q}}{T} + \dot{S}_{gen}
}\end{equation}

where:

\begin{itemize}
\tightlist
\item
  \(\dot{Q}/T\) accounts for entropy transfer with heat
\item
  \(\dot{S}_{gen}\) is the rate of entropy generation
\end{itemize}

Multiply by \(T\) to get:

\[
T\frac{dS}{dt} = \dot{Q} + T\dot{S}_{gen}
\]

Solve for \(\dot{Q}\):

\[
\dot{Q} = T\frac{dS}{dt} - T\dot{S}_{gen}
\]

Substitute this into the energy balance Eq.~\ref{eq-energy-balance}:

\[
\frac{dU}{dt} = T\frac{dS}{dt} - T\dot{S}_{gen} - P\frac{dV}{dt}
\]

Rearrange:

\[
\frac{dU}{dt} + P\frac{dV}{dt} - T\frac{dS}{dt} = -T\dot{S}_{gen}
\]

At constant (P) and constant (T), we may move (P) and (T) inside the
derivative:

\[
\frac{d}{dt}(U+PV-TS) = -T\dot{S}_{gen}
\]

But by definition,

\[
G = U + PV - TS
\]

Therefore,

\begin{equation}\phantomsection\label{eq-Gibbs-entropy-generation}{
\frac{dG}{dt} = -T\dot{S}_{gen}
}\end{equation}

\begin{center}\rule{0.5\linewidth}{0.5pt}\end{center}

\subsection{Key Result}\label{key-result}

\begin{tcolorbox}[enhanced jigsaw, opacitybacktitle=0.6, leftrule=.75mm, left=2mm, toprule=.15mm, toptitle=1mm, colback=white, colbacktitle=quarto-callout-important-color!10!white, bottomrule=.15mm, rightrule=.15mm, colframe=quarto-callout-important-color-frame, arc=.35mm, coltitle=black, breakable, bottomtitle=1mm, titlerule=0mm, title=\textcolor{quarto-callout-important-color}{\faExclamation}\hspace{0.5em}{Key Result: Gibbs Energy and Entropy Generation}, opacityback=0]

At constant \(T\) and \(P\), Eq.~\ref{eq-Gibbs-entropy-generation} shows
that the rate of change of Gibbs energy is directly related to the rate
of entropy generation:

\[
\frac{dG}{dt} = -T\dot{S}_{gen}
\]

Since the second law requires

\[
\dot{S}_{gen} \ge 0
\]

it follows that

\[
\frac{dG}{dt} \le 0
\]

for any spontaneous process.

\end{tcolorbox}

\begin{center}\rule{0.5\linewidth}{0.5pt}\end{center}

\section{\texorpdfstring{Interpretation: Why \(G\) is Minimized at
Equilibrium}{Interpretation: Why G is Minimized at Equilibrium}}\label{interpretation-why-g-is-minimized-at-equilibrium}

This result gives immediate physical insight.

\subsection{Spontaneous process}\label{spontaneous-process}

If the process is spontaneous,

\[
\dot{S}_{gen} > 0
\]

so

\[
\frac{dG}{dt} < 0
\]

That is, \textbf{Gibbs energy decreases with time} for a spontaneous
process.

\subsection{Equilibrium}\label{equilibrium}

At equilibrium,

\[
\dot{S}_{gen} = 0
\]

so

\[
\frac{dG}{dt} = 0
\]

Thus, at constant \(T\) and \(P\), a spontaneous process moves in the
direction of decreasing Gibbs energy, and equilibrium corresponds to a
minimum in \(G\):

\begin{quote}
A spontaneous process moves in the direction of decreasing Gibbs energy,
and equilibrium corresponds to a minimum in (G).
\end{quote}

\begin{center}\rule{0.5\linewidth}{0.5pt}\end{center}

\section{Phase Stability and Direction of Phase
Change}\label{phase-stability-and-direction-of-phase-change}

This explains why \(G\) tells us which phase is stable.

At a given \(T\) and \(P\):

\begin{itemize}
\tightlist
\item
  if \(\underline{G}^v > \underline{G}^l\), vapor should spontaneously
  condense to liquid,
\item
  if \(\underline{G}^l > \underline{G}^v\), liquid should spontaneously
  vaporize,
\item
  if \(\underline{G}^v = \underline{G}^l\), the two phases are in
  equilibrium.
\end{itemize}

So \(G\) tells us the direction a process wants to go.

\begin{quote}
The most stable phase at equilibrium is the one with the lowest molar
Gibbs energy.
\end{quote}

Recall when we were solving cubic equations of state, we had to
determine which root corresponded to the stable phase. We needed to
compare the pressure to the experimental vapor pressure to choose the
proper root, but that assumed the cubic eequation of state captured the
experimental vapor pressure. Now we see that the root with the lowest
molar Gibbs energy is the stable phase. We will come back to this later
when we discuss the concept of fugacity.

\begin{center}\rule{0.5\linewidth}{0.5pt}\end{center}

\subsection{Foreshadowing Mixtures}\label{foreshadowing-mixtures}

For a pure fluid,

\[
\mu = \left(\frac{\partial G}{\partial N}\right)_{T,P}
\qquad \text{and} \qquad
G = N\underline{G}
\]

For a mixture, each species has its own chemical potential:

\[
\mu_i \equiv \left(\frac{\partial G}{\partial N_i}\right)_{T,P,N_{j\ne i}}
\]

That is, when differentiating with respect to \(N_i\), we hold:

\begin{itemize}
\tightlist
\item
  \(T\) constant,
\item
  \(P\) constant,
\item
  all other mole numbers \(N_j\) constant.
\end{itemize}

This will be extremely important when we discuss mixtures and fugacity.

\begin{center}\rule{0.5\linewidth}{0.5pt}\end{center}

\section{Derivation of the Clapeyron
Equation}\label{derivation-of-the-clapeyron-equation}

Now we return to phase equilibrium and ask:

\begin{quote}
How does the equilibrium pressure change with temperature along a phase
boundary?
\end{quote}

Consider a pure fluid on a (P)-(T) phase diagram.

\subsection{Figure placeholder}\label{figure-placeholder}

\textbf{{[}Insert figure here: a (P)-(T) phase diagram showing a
coexistence line between two phases (\alpha) and (\beta), with a point
on the line and small changes (dT) and (dP){]}}

Suppose the system is initially at a point on the coexistence line. If
we change the temperature by (dT), then the pressure must also change by
(dP) in order to remain on the phase boundary.

For a pure fluid,

\[
d\underline{G} = \underline{V}\,dP - \underline{S}\,dT
\]

At equilibrium between phases (\alpha) and (\beta),

\[
\underline{G}^\alpha = \underline{G}^\beta
\]

and if we remain on the coexistence curve,

\[
d\underline{G}^\alpha = d\underline{G}^\beta
\]

So:

\[
\underline{V}^\alpha dP - \underline{S}^\alpha dT
=
\underline{V}^\beta dP - \underline{S}^\beta dT
\]

Rearranging:

\[
(\underline{V}^\beta - \underline{V}^\alpha)dP
=
(\underline{S}^\beta - \underline{S}^\alpha)dT
\]

Therefore,

\[
\left(\frac{dP}{dT}\right)_{\alpha-\beta}
=
\frac{\underline{S}^\beta - \underline{S}^\alpha}{\underline{V}^\beta - \underline{V}^\alpha}
\]

or more compactly,

\[
\left(\frac{dP}{dT}\right)_{\alpha-\beta}
=
\frac{\Delta \underline{S}^{\alpha\to\beta}}{\Delta \underline{V}^{\alpha\to\beta}}
\]

This is one form of the \textbf{Clapeyron equation}.

\begin{center}\rule{0.5\linewidth}{0.5pt}\end{center}

\section{Clapeyron Equation in Terms of
Enthalpy}\label{clapeyron-equation-in-terms-of-enthalpy}

Using the fact that for a reversible phase change at equilibrium,

\[
\Delta \underline{S}^{\alpha\to\beta}
=
\frac{\Delta \underline{H}^{\alpha\to\beta}}{T}
\]

we obtain

\[
\left(\frac{dP}{dT}\right)_{\alpha-\beta}
=
\frac{\Delta \underline{H}^{\alpha\to\beta}}{T\,\Delta \underline{V}^{\alpha\to\beta}}
\]

This is the most common form of the \textbf{Clapeyron equation}:

\[
\boxed{
\left(\frac{dP}{dT}\right)_{\alpha-\beta}
=
\frac{\Delta \underline{H}^{\alpha\to\beta}}{T\,\Delta \underline{V}^{\alpha\to\beta}}
}
\]

It is an \textbf{exact} equation and applies to \textbf{any two phases
in equilibrium}.

\begin{center}\rule{0.5\linewidth}{0.5pt}\end{center}

\section{Physical Implications of the Clapeyron
Equation}\label{physical-implications-of-the-clapeyron-equation}

The sign of the slope depends on the signs of (\Delta \underline{H}) and
(\Delta \underline{V}).

\subsection{Vapor-liquid equilibrium}\label{vapor-liquid-equilibrium}

For vaporization,

\begin{itemize}
\tightlist
\item
  (\Delta \underline{H}\_\{vap\} \textgreater{} 0),
\item
  (\Delta \underline{V}\_\{vap\} \textgreater{} 0),
\end{itemize}

so

\[
\left(\frac{dP}{dT}\right)_{VLE} > 0
\]

Thus the vapor-liquid line has a positive slope.

\subsection{Solid-liquid equilibrium}\label{solid-liquid-equilibrium}

For melting,

\begin{itemize}
\tightlist
\item
  (\Delta \underline{H}\_\{fus\} \textgreater{} 0) almost always,
\item
  (\Delta \underline{V}\_\{fus\}) is usually positive, but not always.
\end{itemize}

For most substances,

\[
\underline{V}^l > \underline{V}^s
\]

so

\[
\left(\frac{dP}{dT}\right)_{SLE} > 0
\]

\subsection{Why water is unusual}\label{why-water-is-unusual}

Water is unusual because ice is less dense than liquid water, so

\[
\underline{V}^s > \underline{V}^l
\]

Thus, for water,

\[
\Delta \underline{V}_{fus} = \underline{V}^l - \underline{V}^s < 0
\]

and therefore

\[
\left(\frac{dP}{dT}\right)_{SLE} < 0
\]

So the solid-liquid line for water has a \textbf{negative slope}.

\subsection{Figure placeholder}\label{figure-placeholder-1}

\textbf{{[}Insert figure here: phase diagram of water highlighting the
negative slope of the solid-liquid coexistence line{]}}

Other substances with this behavior include bismuth, gallium, germanium,
silicon, and arsenic.

\begin{center}\rule{0.5\linewidth}{0.5pt}\end{center}

\section{Derivation of the Clausius--Clapeyron
Equation}\label{derivation-of-the-clausiusclapeyron-equation}

Now specialize the Clapeyron equation to \textbf{vapor-liquid
equilibrium}.

Start with

\[
\left(\frac{dP}{dT}\right)_{VLE}
=
\frac{\Delta \underline{H}_{vap}}{T(\underline{V}^v-\underline{V}^l)}
\]

To simplify this, make three assumptions.

\subsection{Assumption 1: vapor volume dominates liquid
volume}\label{assumption-1-vapor-volume-dominates-liquid-volume}

Assume

\[
\underline{V}^v \gg \underline{V}^l
\]

so

\[
\underline{V}^v - \underline{V}^l \approx \underline{V}^v
\]

This is usually a good approximation except near the critical point,
where the liquid and vapor densities become similar.

\subsection{Assumption 2: vapor behaves
ideally}\label{assumption-2-vapor-behaves-ideally}

Assume the vapor is an ideal gas:

\[
\underline{V}^v = \frac{RT}{P}
\]

Combining the first two assumptions:

\[
\underline{V}^v - \underline{V}^l \approx \underline{V}^v \approx \frac{RT}{P}
\]

Substitute into the Clapeyron equation:

\[
\frac{dP^{sat}}{dT}
\approx
\frac{\Delta \underline{H}_{vap}}{T(RT/P^{sat})}
\]

so

\[
\frac{dP^{sat}}{dT}
\approx
\frac{P^{sat}\Delta \underline{H}_{vap}}{RT^2}
\]

Rearrange:

\[
\frac{dP^{sat}}{P^{sat}}
\approx
\frac{\Delta \underline{H}_{vap}}{R}\frac{dT}{T^2}
\]

or equivalently,

\[
d\ln P^{sat}
=
-\frac{\Delta \underline{H}_{vap}}{R}\,d\left(\frac{1}{T}\right)
\]

\subsection{\texorpdfstring{Assumption 3:
(\Delta \underline{H}\_\{vap\}) is
constant}{Assumption 3: (\_\{vap\}) is constant}}\label{assumption-3-_vap-is-constant}

Assume (\Delta \underline{H}\_\{vap\}) does not vary strongly with
temperature over the range of interest. Then integration gives:

\[
\ln\left(\frac{P_2^{sat}}{P_1^{sat}}\right)
=
-\frac{\Delta \underline{H}_{vap}}{R}
\left(\frac{1}{T_2}-\frac{1}{T_1}\right)
\]

This is the \textbf{Clausius--Clapeyron equation}.

\begin{center}\rule{0.5\linewidth}{0.5pt}\end{center}

\section{Clausius--Clapeyron Equation}\label{clausiusclapeyron-equation}

The Clausius--Clapeyron equation may be written as

\[
\boxed{
\ln\left(\frac{P_2^{sat}}{P_1^{sat}}\right)
=
-\frac{\Delta \underline{H}_{vap}}{R}
\left(\frac{1}{T_2}-\frac{1}{T_1}\right)
}
\]

or, relative to an arbitrary constant,

\[
\ln P^{sat} = -\frac{\Delta \underline{H}_{vap}}{R}\frac{1}{T} + C
\]

This is an \textbf{approximate} equation, because it relies on the three
assumptions above, but it is often very useful and reasonably accurate
over moderate temperature ranges.

\begin{center}\rule{0.5\linewidth}{0.5pt}\end{center}

\section{Interpretation of
Clausius--Clapeyron}\label{interpretation-of-clausiusclapeyron}

This equation shows that a plot of

\[
\ln P^{sat} \quad \text{vs.} \quad \frac{1}{T}
\]

should be approximately linear, with slope

\[
-\frac{\Delta \underline{H}_{vap}}{R}
\]

Therefore, vapor pressure data can be used to estimate the enthalpy of
vaporization.

\subsection{Figure placeholder}\label{figure-placeholder-2}

\textbf{{[}Insert figure here: straight-line plot of (\ln P\^{}\{sat\})
versus (1/T), with slope labeled as
(-\Delta \underline{H}\_\{vap\}/R){]}}

This also helps explain why the Antoine equation has the form it does.

\begin{center}\rule{0.5\linewidth}{0.5pt}\end{center}

\section{Example: Why Is Ice
Slippery?}\label{example-why-is-ice-slippery}

A classic hypothesis is that the pressure exerted by an ice skate melts
the ice underneath the blade, creating a thin liquid water layer that
reduces friction.

Let us test whether this is plausible using the Clapeyron equation.

Suppose ice is initially at (-5\^{}\circ\text{C}) and 1 bar. How much
pressure would be needed to make it melt at (-5\^{}\circ\text{C})?

We cannot use the Clausius--Clapeyron equation here because this is
\textbf{solid-liquid equilibrium}, not vapor-liquid equilibrium.

Instead use the exact Clapeyron equation:

\[
\left(\frac{dP}{dT}\right)_{SLE}
=
\frac{\Delta \underline{H}_{fus}}{T\,\Delta \underline{V}_{fus}}
\]

Assume (\Delta \underline{H}\emph{\{fus\}) and
(\Delta \underline{V}}\{fus\}) are approximately constant over the small
temperature range.

Using data near (0\^{}\circ\text{C}):

\begin{itemize}
\tightlist
\item
  (\rho\^{}l \approx 1000~\text{kg/m}\^{}3)
\item
  (\rho\^{}s \approx 917~\text{kg/m}\^{}3)
\item
  (\Delta \underline{H}\_\{fus\} \approx 6010~\text{J/mol})
\end{itemize}

The molar volumes are approximately

\[
\underline{V}^l \approx 1.80\times10^{-5}\ \text{m}^3/\text{mol}
\]

\[
\underline{V}^s \approx 1.97\times10^{-5}\ \text{m}^3/\text{mol}
\]

so

\[
\Delta \underline{V}_{fus} = \underline{V}^l-\underline{V}^s < 0
\]

Integrating the Clapeyron equation approximately gives a required
pressure on the order of tens of bar, roughly

\[
P_2 \approx 66\ \text{bar}
\]

This is much larger than the pressure exerted by a typical skate blade,
which is only a few bar.

Therefore:

\begin{quote}
Pressure melting alone does \textbf{not} explain why ice is slippery.
\end{quote}

A better explanation is that the ice surface is not a perfect bulk
solid; a thin disordered or quasi-liquid interfacial layer can exist,
and frictional heating may also contribute.

\subsection{Figure placeholder}\label{figure-placeholder-3}

\textbf{{[}Insert figure here: sketch of skate blade on ice and
comparison of estimated skate pressure to pressure required for melting
at (-5\^{}\circ\text{C}){]}}

\begin{center}\rule{0.5\linewidth}{0.5pt}\end{center}

\section{Summary}\label{summary}

\begin{itemize}
\item
  Chemical potential is the conjugate variable to (N) and can be written
  as: \[
  \mu =
  \left(\frac{\partial U}{\partial N}\right)_{S,V}
  =
  \left(\frac{\partial H}{\partial N}\right)_{S,P}
  =
  \left(\frac{\partial A}{\partial N}\right)_{T,V}
  =
  \left(\frac{\partial G}{\partial N}\right)_{T,P}
  \]
\item
  The general fundamental relations are: \[
  dU = T\,dS - P\,dV + \mu\,dN
  \] \[
  dH = T\,dS + V\,dP + \mu\,dN
  \] \[
  dA = -S\,dT - P\,dV + \mu\,dN
  \] \[
  dG = -S\,dT + V\,dP + \mu\,dN
  \]
\item
  For a closed system at constant (T) and (P), \[
  \frac{dG}{dt} = -T\dot{S}_{gen}
  \]
\item
  Therefore:

  \begin{itemize}
  \tightlist
  \item
    spontaneous processes decrease (G),
  \item
    equilibrium corresponds to a minimum in (G).
  \end{itemize}
\item
  The exact Clapeyron equation is: \[
  \left(\frac{dP}{dT}\right)_{\alpha-\beta}
  =
  \frac{\Delta \underline{H}^{\alpha\to\beta}}{T\,\Delta \underline{V}^{\alpha\to\beta}}
  \]
\item
  For vapor-liquid equilibrium, under additional assumptions, this
  becomes the Clausius--Clapeyron equation: \[
  \ln\left(\frac{P_2^{sat}}{P_1^{sat}}\right)
  =
  -\frac{\Delta \underline{H}_{vap}}{R}
  \left(\frac{1}{T_2}-\frac{1}{T_1}\right)
  \]
\end{itemize}

\begin{center}\rule{0.5\linewidth}{0.5pt}\end{center}

\section{Looking Ahead}\label{looking-ahead}

Next we will use these ideas to develop practical tools for real-fluid
phase equilibrium calculations, including:

\begin{itemize}
\tightlist
\item
  fugacity,
\item
  Antoine equation,
\item
  and other useful correlations for vapor pressure.
\end{itemize}




\end{document}
